\documentclass[]{../../../../tex/classes/styledArticle}
\usepackage{../../../../tex/packages/hebrewSupport}


\author{שחר פרץ}
\title{\textit{היסטוריה 51}}
\begin{document}
	\maketitle
	בעשור הראשון, היו כאן 650K יהודים שלקטו מיליון (עד 1960). מאיפה הגיעו העולים? מכל מקום. 
	\begin{itemize}
		\item אירופה (שורדי שואה) – 420 אלף עולים. מרומניה, פולין, בולגריה, הונגריה, ברה''מ, צ'כוסלובקיה, גרמניה, אוסטריה, יוגוסלביה, צרפת, איטליה, הולנגד, בלגיה ותורכיה. ממזרח אירופה (ברה''מ) השערים נסגרו בשלב מסויים ולא התאפשרה יותר עלייה, ועד נפילת מסך הברזל לא הייתה עליה מברה''מ. 
		\item מאסיה 275K – תורכיה, סוריה, לבנון, עיראק, איראן, הודו ואפגניסטן. 
		\item מאפריקה 230K – מרוקו, אלג'יריה, תוניס, לוב, מצרים, סודן, דרום אפריקה. 
	\end{itemize}
	
	כל הסיפור הזה התאפשר מחוק השבות – כל אזרח יהודי יכול לעלות ולקבל אזרחות. ומיהו יהודי? ההגדרה יותר מרחיבה מההגדרה האורתודוקסית – זה יכול להיות גם לפי האב (ולא לפי האם) ובגיוס רפורמי (נכון להיום, הממשלה מנסה לשנות את זה). זה לא אומר שמבחינת הרבנות תהיה יהודי. בצבא הייתה תקופה של עלייה גדולה מברה''מ, שע''פ היהדות הם לא היו יהודים, ובצבא היה (ועדיין פועל) פרויקט גיור מסודר של אותם האנשים (בהתנדבות). כשמדובר על עליה, מדובר על עליה יהודית על סמך חוק השבות, לא לפי הרבנות. 
	
	הגורמים לעלייה – מה שאנחנו מכירים. 
	\begin{itemize}
		\item הרבה מאוד עולים הם ניצולי שואה שבקשו לברוח מאירופה, או לא יכלו להישאר באירופה (הבתים היו תפוסים, פוגרומים בפולין, וכו'). 
		\item החמרה ביחס של הערבים כלפי יהודים בארצות האסלאם עם הקמת מדינת ישראל, אילצה יהודים רבים לעזוב לא''י. 
		\item מאבק של ארצות האסלאם לעצמאות בשלטון הקולוניאלי והתחזקות הפאן־ערביות והפאן־איסלאמיות. 
		\item יחס עוין כלפי יהודי במזרח אירופה לאחר מרידות ותסיסות אנטי־קומוניסטיות (היהודים במזרח אירופה בחלקם הגדול היו קומוניסטים או נתפסו כחלק מהגוש הקומוניסטי). 
	\end{itemize}
	חלק מיהודים שעוזבים את מדינת האסלאם לא מגיעים לא''י. בקהילת לוב, היו מן היהודים שתפסו את עצמם כאיטלקים ונטשו לאיטליה, ולא לא''י – ויש קהילה עד היום לא קטנה של עולי לוב ברומא. שימו לב שלא כל הזמן היהודים ממזרח אירופה יכלו לעלות – קראו על מסך הברזל. 
	
	דפוסי/מאפייני העליה: 
	\begin{itemize}
		\item עליית הצלה – חיסול גלויות. במקרים מסויימים, מנהיגי המדינה חשבו שנשקפת סכנה קיומית לחייהם של היהודים בארץ מסויימת ומצבם רע מאוד, ולכן צריך להעלות את כל הקהילה היודית בזמן קצר לישראל. לדוגמה בעיראק, לוב, ואיראן (מאוחר יותר). המאפיין של עלייה זו היה הדאגה לביטחונם של היהודים. מדינת ישראל העלתה קהילות שלמות לארץ. דיברנו על זה כשדיברנו על גודל יהודי ארצות האסלאם. בחלק מהמדינות זה בתיאום על השלטונות (בתנאים מסויימים, כמו שישארו את רכושם). 
		\item עלייה המונית – שערי א''י פתוחים, מי שבא נכנס. 
		\item עליית (הצלה) חשאית – עלייה לא לגאלית, בכלל סכנה קיומית ליהודים בארצות מוצאם, למשל מסוריה, לבנון, ומצריים בין 57 ל־61. מדובר על עלייה בלתית לגאלית, לרוב עליית הצלה עברה דרך מדינת צד שלישי. במקרה של לבנון וסוריה וכו', לא היה צורך במטוסים וכו', כי יש גבול. 
		\item עלייה מבוקרת/סלקטיבית – אם לא נשקפה סכנה מידית לקהילה היהודית, מנהיגי המדינה העדיפו לשמור על האינטרסים של המדינה. היה חשש שהמדינה לא תצליח לקלוט כמות גדולה של עולים, ולכן, התחילו במדיניות סלקציה. 
	\end{itemize}
	
	בתחילת שנות השבעים אפשרו לאנשים משטחי רוסיה (לא ממדינות חבר העמים) לעלות לארץ לזמן קצר. 
	
	למה עשו עלייה מבוקרת? בעלייה המונית / עליית הצלה, עולים המון עולים שלא יכולים לעזור למדינה ולתחזק את עצמם (מבוגרים וכו') ויעלו את הצעירים שיכולים לעבוד. לדוגמה, במרוקו העלייה התחילה כעליה סלקטיבית. אנשים שתלויים באחרים, ישארו במדינות מוצאם, ואחת שהמדינה תתפתח דיו הם יעלו לארץ. בן־גוריון בצורה מאוד בולטת לא סבל עלייה סלקטיבית. הוא היה בעד עלייה המונית. 
	
	\subsection*{תהליך הקליטה}
	היו קשיים מצד המדינה בעלייה. 
	\begin{itemize}
		\item היקף העלייה
		\item ארגון העלייה – מגיעים המון אנשים, צריך לדאוג להם למגורים, אוכל, תנאי מחיה, פרנסה, וכו'. בשנה הראשונה, המדינה בכלל במלחמה. אחד מהאינטרסים לחתום על הסכמי שביתת הנשק הייתה העלייה, כי רוב התקציב הקטן ממילא של המדינה הלך לצרכים צבאיים. בהתחלה לא היו מוסדות לקליטת עלייה, הכל היה בהתהוות. 
		\item הקשיים וההשתלבות החברתית – רוב העולים היו אשכנזים, היו פערי תרבות גדולים ביחס לעולים. 
	\end{itemize}
	
	גם לעולים יש קשיים
	\begin{multicols}{2}
		\begin{itemize}
			\item מעבר מארץ לידה לארץ חדשה
			\item קשיי שפה, בעיקר בקרב המבוגרים
			\item מפגש עם תרבות אחרת (החברה הקולטת היא ממזרח אירופה ברובה)
			\item עליית משפחות – ילדים, מבוגרים וניצולי שואה
			\item קושי במציאת תעסוקה
			\item ערעור חיי המשפחה – פרטים בהמשך
			\item תחושת קיפוח ואפליה (חברת הנקלטים לעומת הקולטים)
		\end{itemize}
	\end{multicols}
	''הכי קל ללמד את זה בנעל''ה``
	
	אנשים רבים היו מנהיגי קהילה, עם קשרים לאימפריה, משרתים בבית, וכיו''ב, והם עולים לפחונים. רוב היהודים עלו בלי רכושם. תחשבו מה זה עושה למבנה המשפחה. במדינות האסלאם זה היה יותר גרוע – שם, המבנה המשפחתי היה פטריארכלי, וחשבו על האבא ''שפעם היה וואו, והיום...``. מצב שבו אבא נהיה תלוי בבן שלו, שיצא ממבנה משפחתי פטריארכלי, היה די נורא. סתם דוגמה של שרית – ביום הורים, ילד מגיע, ההורים בכלל לא מבינים את שרית כי הם לא דוברים עברית. אין להם מושג מה הילד מרגיש, הם לא מכירים את התרבות. הילדים מתערים מהר, לומדים את השפה והתרבות, והוריהם מתקשים להתקדם. אנשים שהגיעו עם מקצועות (פחות ממדינות מוסלמיות) לדוגמה רופאים, מהנדסים וכו', מצאו עצמם מנקים. תחשבו שעולים לארץ הקודש זבת חלב ודבש ומגיעים למעברה, ''שם יפה למחנה פליטים, בתנאים של מחנה פליטים``. ילדים נכנסו למוסדות לימוד/צבא, דהיינו, כור היתוך על סטרואידים, 
	
	אחד מהדברים שמדברים עליהם כשמדברים על בעיית הפליטים הערבים, היא שהמדינות יכלו לקלוט את הפליטים – כמו שישראל קלטה את הפליטים היהודים, לא קלטה אותם, והם חיים במחנות פליטים עד היום (עם יוצא דופן של ירדן, שם יש מעט מאוד מחנות פליטים שנשארו). מסיבות פוליטיות. 
	
	מי שלמד על ה־new deal, גם במקרה של מדינת ישראל, המדינה שלמה משכורת מינימלית בשביל שאנשים יתפרנסו – סלילת כבישים, ייעור, וכו'. 
	
	<ראינו קטע מלול בשיעור שעבר>
	
	<רואים קטע מסלאח שבתי, ''שאגב צולם בכפר``> ''אריק איינשטיין. הוא היה חתיך. לא לזלזל``
	
	\begin{itemize}
		\item התמודדות עם היקף העלייה – לדוגמה באמצעות עלייה סלקטיבית
		\item התמודדות עם הדיור: 
		\begin{enumerate}
			\item ראשית כל, לקלוט אנשים בבתים של ערבים שננטשו במלחמה. האנשים שהגיעו לארץ ראשונים, נכנסו לבתים של ערבים שהפכו להיות בבעלותם. באופן יחסי לכמות העולים, לא היו המון בתים כאלו. 
			\item מחנות עולים – מחנות צבאיים בריטיים, שהבריטים פינו, ששמו שם את העולים. במחנות האלו, התלות של העולים היא תלות מלאה לחלוטין במדינה – הם לא יוצאים לעבוד, ויש שם את כל מה שהם צריכים (חדרי אוכל, מרפאות וכיו''ב). 
			\item המעברות – יש כל מני סוגים, פחונים, צריפים (מעץ), אוהלים, בדונים ועוד. המעברות בד''כ היו צמודות ליישובים קיימים, כדי שהאנשים שחיים במעברות יוכלו לצאת ולהתפרנס. <עמירם פותח חזית מול שרית> במעברות היו שירותים מרכזיים, כמו לשכת תעסוקה, מרפאה, בתי עבודה סוציאלית, כיתות לימוד, וכו'. לדוגמה, ליד אבן יהודה הייתה מעברה ענקית. 
			\item מושבי עולים ועיירות פיתוח, שהוקמו כרצועת בטחון. (דניאלה ר אומרת ''הרצליה`` [פיתוח]. צוחקים.) ''אני בטוחה שהאנשים באופקים היו שמחים לגור בהרצליה פיתוח``. המקומות האלו הוקמו בגבולות, ונועדו לפתח את החקלאות (לדוגמה, אופקים, קריית שמונה). זו התיישבות בטחונית. ''במקום שבו עוברת המחרשה, עובר הגבול`` – אי אפשר לשים חייל על כל מאה מטר, אז באותה התקופה, ההתיישבות היא זו ששמרה בין היתר על גבולות המדינה. 
		\end{enumerate}
		למושבי העולים הביאו עולים חדשים ''שבינם לבין חקלאות לא היה שום קשר, וקצת הכריחו אותם לעבוד בחקלאות``. הרבה פעמים הביאו לשם אוכלוסיות הומוגניות – אז היה יישוב של עולים מלוב, של עולים מעיראק, וכו'. 
		\item התמודדות עם הכלכלה והתעסוקה – 
		\begin{itemize}
			\item מגביות (איסוף תרומות) מיהודי התפוצות. המון תרומות מיהודי בחו''ל. 
			\item משטר צנע – היער המשאבים הכלכליים אילץ את הממשלה להכריז על תוכנית חירום, במסגרתה הוטל פיקוח על הצריכה והמחירים של המזון. הייתה הגבלה על המזון פר אדם. בתקופה הזו נולדו פתיתים, ''אורז בן־גוריון``. ''יש הבדל בין פסטה לפתיתים``, שרית מסבירה לליה. גלי רימר מתחילה לדבר על פסטינה בניסיון לשכנע את שרית שלא אנחנו המצאנו את זה. שבוע הבא יובל יביא פתיתים, גלי תביא פסטינה, ונעשה השוואה. 
			\item כספי השילומים – פנייה שיזם בן גוריון לגרמניה בגרישה להחזיר את כספי היהודים שנלקחו מהם בשואה. היה כאן ויכוח שלם, כמעט קרע במדינה, בנושא הזה. הפיצויים האלו התחלקו לשני חלקים – חלקם ניתן למדינת ישראל, וחלקם שורדי השואה מקבלים באופן אישי. 
			\item עבודות יזומות – מצד אחד, עוזרות לפתח את המדינה, ומצד שני, נותנת פרנסה מינימלית ובסיסית לאנשים שאין להם עבודה. 
		\end{itemize}
		\item התמודדות עם בעיות הביטחון: 
		\begin{itemize}
			\item עיירות פיתוח שהוקמו בגבולות
			\item חוק גיוס חובה
		\end{itemize}
		\item חברה וחינוך: 
		\begin{itemize}
			\item חוק חינוך חובה (1949, ידוע בכינויו ''חחח``) המחייב את התלמידים מגיל 5-14 ללמוד בבית הספר, מחייב את המדינה לתת לו חינוך, ומחייב את ההורים לשלוח אותנו לביה''ס. ''באופן תיאורטי, לדעתי זה לעולם לא קרה, יכולים לשלוח את ההורים שלכם לכלא``. 
			\item הנחלת הלשון העברית – מבצע שהחל ב־1954 ובמהלכו 60K אנשים למדו עברית. 
		\end{itemize}
		\item דרכי התמודדות עם הקליטה התרבותית והחברתית: כור ההיתוך. דניאלה רחמנפור הייתה בטוחה שמדובר על מפעל הטקסטיל בדימונה. בכור היתוך מייצרים טקסטיל. בוהים בציור המפורסם בו אנשים שמחים ומחייכים נכנסים לתוך כור ההיתוך ויוצא שרוליק. התפיסה – כולם צריכים להיות שרוליק. מה הבעיה? שרוליק אשכנזי. הבסיס לרעיון הצבר הוא אשכנזי, כי הרוב כאן היה אשכנזי והם אלו שעצבו את התפיסה. הכוונה הייתה טובה – יש כאן אנשים מתרבויות שונות, עם הרגלים שונים, ורוצים להפוך אותם לחברה אחת. המטרות: 
		\begin{itemize}
			\item לצמצם את הניכור
			\item ליצור דמות חדשה ומודרנית של ישראלי
			\item לעצב תרבות חדשה שאינה דומה לאף מרכיב קודם
		\end{itemize}
		כל אחד היה אמור להשיל מעליו את התרבות ממנו הוא בא, ולהפוך להיות צבר. איך עושים את זה? חינוך וצה''ל. מפמפמים ערכים לאומיים, תכנית ''אינטגרציה בחינוך`` (תוכנית בשנות ה־70 וה־80 בה הוקמו חטיבות ביניים עם כיתות הטרוגניות). בזמן כור ההיתוך, אם אמא דיברה איטלקית ולא עברית, זה היה מביך. המדיניות כמובן חלה על כולם – כולל יהודי שורדי השואה שנתקעו במעברות. מדברים פחות על האשכנזים כי הם לכאורה משוייכים לצד של השלטון, והתפיסה של הצבר הייתה יותר קרובה אליהם תרבותית. אבל היחס היה באותה המידה למי שהגיע מאירופה ומי שהגיע ממדינות מוסלמיות. ''אוי ואבוי אם דיברת יידיש, לדינו, איטלקית, רוסית ומרוקאית`` – זה נוגד את התפיסה. צה''ל עד היום הוא כור היתוך מטורף, ומי שהכי מצליח לשלב אנשים בחברה הישראלית, הוא צה''ל – כשאתה נכנס למערכת, זה לא מעניין אף אחד מאיפה אתה. ברגע שנכנסת לצה''ל, כולם מתחילים מאותו המקום. הבעיה היא לפני צה''ל – כי נגיד ל־8200, כנראה זה פחות נפוץ משדרות מאשר בכפר. זה קשור בחינוך לפני, לא בצה''ל. 
		
		אירועי וואדי סליב – תקראו. עוד מערכונים של סלאח שבתי פספסנו. 
	\end{itemize}
	
	חוץ מבעיית הפליטים והסכמי שביתת הנשק, ומלחמת ששת הימים, למדנו הכל. 
	
	
\end{document}